% ============================================================
% Chapter 2: State of the Art
%   Section: Benchmarks for Fact Verification
% Included from chapters/02-state-of-the-art/state-of-the-art.tex
% ============================================================

\section{Benchmarks for Fact Verification}
\label{sec:sota-benchmarks}

Table~\ref{tab:sota-benchmarks} summarises the benchmarks most relevant
to this work. They differ in four respects that decide whether they can
evaluate an open-web verifier of constructive news: where the claims
come from, where the evidence comes from, the label scale, and the
languages.

\begin{table}[ht]
\centering
\small
\begin{tabularx}{\linewidth}{@{}lrXXl@{}}
\toprule
\textbf{Benchmark} & \textbf{Claims} & \textbf{Claims from} &
\textbf{Evidence} & \textbf{Labels} \\
\midrule
FEVER \cite{thorne2018fever} & 185,445 & Altered Wikipedia sentences
(English) & Wikipedia sentences (closed) & 3 \\
LIAR \cite{wang2017liar} & 12.8k & PolitiFact statements (English) &
None; claim and metadata only & 6 \\
MultiFC \cite{augenstein2019multifc} & 34,918 & 26 fact-checking
websites (English) & Search-engine snippets & per site \\
SciFact \cite{wadden2020scifact} & 1.4k & Rewritten citation sentences
(English) & Scientific abstracts (closed) & 3 \\
X-Fact \cite{gupta2021xfact} & 31,189 & Fact-checkers, 25 languages &
Search-engine snippets & 7 \\
AVeriTeC \cite{schlichtkrull2023averitec} & 4,568 & Real-world claims
checked by fact-checkers (English) & Web documents from before the
claim, with question--answer annotations & 4 \\
\bottomrule
\end{tabularx}
\caption{Benchmarks for fact verification relevant to this work.}
\label{tab:sota-benchmarks}
\end{table}

FEVER and SciFact are closed: their evidence is a fixed corpus, which
makes them precise and unlike the task this system performs. LIAR
gives no evidence at all, so it measures what can be guessed from a
claim and its speaker. MultiFC and X-Fact are open, real-world and, in
X-Fact's case, multilingual, with Spanish among the languages of its
test sets; their evidence is search snippets collected when each
dataset was built. AVeriTeC is the closest to how a fact-checker works
--- web evidence, dated before the claim --- but is English only.

Two observations follow for this work. First, every benchmark with
real-world claims takes them from fact-checking organisations, which
choose claims because someone disputed them; most are political, and
none is drawn from constructive news. Second, the benchmarks report a
verdict's quality, not the effort it saves, which is what RQ4 asks.
The choice made here --- X-Fact as an external, open-retrieval
benchmark and a custom validation set for the system's own domain ---
and the reasons for it are in Section~\ref{sec:dataset-selection}.
